\documentclass{amsart}
% For dealing with references we use the comment environment
\usepackage{verbatim}
\newenvironment{reference}{\comment}{\endcomment}
%\newenvironment{reference}{}{}
\newenvironment{slogan}{\comment}{\endcomment}
\newenvironment{history}{\comment}{\endcomment}

% For commutative diagrams we use Xy-pic
\usepackage[all]{xy}

% We use 2cell for 2-commutative diagrams.
\xyoption{2cell}
\UseAllTwocells

% We use multicol for the list of chapters between chapters
\usepackage{multicol}

% This is generally recommended for better output
\usepackage{lmodern}
\usepackage[T1]{fontenc}

% For cross-file-references

% Package for hypertext links:
\usepackage{hyperref}

% For any local file, say "hello.tex" you want to link to please

% Theorem environments.
%
\theoremstyle{plain}
\newtheorem{theorem}[subsection]{Theorem}
\newtheorem{proposition}[subsection]{Proposition}
\newtheorem{lemma}[subsection]{Lemma}

\theoremstyle{definition}
\newtheorem{definition}[subsection]{Definition}
\newtheorem{example}[subsection]{Example}
\newtheorem{exercise}[subsection]{Exercise}
\newtheorem{situation}[subsection]{Situation}

\theoremstyle{remark}
\newtheorem{remark}[subsection]{Remark}
\newtheorem{remarks}[subsection]{Remarks}

\numberwithin{equation}{subsection}

% Macros
%
\def\lim{\mathop{\mathrm{lim}}\nolimits}
\def\colim{\mathop{\mathrm{colim}}\nolimits}
\def\Spec{\mathop{\mathrm{Spec}}}
\def\Hom{\mathop{\mathrm{Hom}}\nolimits}
\def\Ext{\mathop{\mathrm{Ext}}\nolimits}
\def\SheafHom{\mathop{\mathcal{H}\!\mathit{om}}\nolimits}
\def\SheafExt{\mathop{\mathcal{E}\!\mathit{xt}}\nolimits}
\def\Sch{\mathit{Sch}}
\def\Mor{\mathop{\mathrm{Mor}}\nolimits}
\def\Ob{\mathop{\mathrm{Ob}}\nolimits}
\def\Sh{\mathop{\mathit{Sh}}\nolimits}
\def\NL{\mathop{N\!L}\nolimits}
\def\CH{\mathop{\mathrm{CH}}\nolimits}
\def\proetale{{pro\text{-}\acute{e}tale}}
\def\etale{{\acute{e}tale}}
\def\QCoh{\mathit{QCoh}}
\def\Ker{\mathop{\mathrm{Ker}}}
\def\Im{\mathop{\mathrm{Im}}}
\def\Coker{\mathop{\mathrm{Coker}}}
\def\Coim{\mathop{\mathrm{Coim}}}

% Boxtimes
%
\DeclareMathSymbol{\boxtimes}{\mathbin}{AMSa}{"02}

%
% Macros for moduli stacks/spaces
%
\def\QCohstack{\mathcal{QC}\!\mathit{oh}}
\def\Cohstack{\mathcal{C}\!\mathit{oh}}
\def\Spacesstack{\mathcal{S}\!\mathit{paces}}
\def\Quotfunctor{\mathrm{Quot}}
\def\Hilbfunctor{\mathrm{Hilb}}
\def\Curvesstack{\mathcal{C}\!\mathit{urves}}
\def\Polarizedstack{\mathcal{P}\!\mathit{olarized}}
\def\Complexesstack{\mathcal{C}\!\mathit{omplexes}}
% \Pic is the operator that assigns to X its picard group, usage \Pic(X)
% \Picardstack_{X/B} denotes the Picard stack of X over B
% \Picardfunctor_{X/B} denotes the Picard functor of X over B
\def\Pic{\mathop{\mathrm{Pic}}\nolimits}
\def\Picardstack{\mathcal{P}\!\mathit{ic}}
\def\Picardfunctor{\mathrm{Pic}}
\def\Deformationcategory{\mathcal{D}\!\mathit{ef}}

\setlength{\emergencystretch}{3em}
\begin{document}
\title[Stacks Project, Tag 06HL]{397 --- Smooth morphisms of prorepresentable functors are power-series extensions}
\maketitle
\noindent Copyright (C) 2005--2025 Johan de Jong. Stacks Project authors. Source: \href{https://stacks.math.columbia.edu/tag/06HL}{Tag 06HL}.

\noindent Editorial difficulty note (not author text). Difficulty H2: The proof lifts a relative cotangent-space isomorphism successively through Artinian quotients to obtain a map into a power-series ring. It then constructs a section and uses complete-local cotangent-space surjectivity twice to prove both maps are inverse isomorphisms. The formal smoothness characterization requires this completed lifting and rigidity construction..

\noindent Editorial provenance note (not author text). History: excerpt from revision \texttt{a0444\allowbreak 6e57e\allowbreak c1fbc\allowbreak 252a8\allowbreak 71afc\allowbreak ec775\allowbreak 2fb28\allowbreak 07b14}, formal-defos.tex, lines 2024--2102. Reference presentation was adapted; original-author context and core support blocks were retained or added. The mathematical wording is unchanged. Licensed under GNU FDL 1.2 or later; no invariant sections or cover texts. See ../licenses/COPYING. Context blocks reproduce author definitions and supporting results; they are not additional counted problems.

\begin{definition}
\label{definition-CLambda}
Let $\Lambda$ be a Noetherian ring and let $\Lambda \to k$ be a finite
ring map where $k$ is a field. We define {\it $\mathcal{C}_\Lambda$} to be
the category with
\begin{enumerate}
\item objects are pairs $(A, \varphi)$ where $A$ is an Artinian local
$\Lambda$-algebra and where $\varphi : A/\mathfrak m_A \to k$ is a
$\Lambda$-algebra isomorphism, and
\item morphisms $f : (B, \psi) \to (A, \varphi)$ are local $\Lambda$-algebra
homomorphisms such that $\varphi \circ (f \bmod \mathfrak m) = \psi$.
\end{enumerate}
We say we are in the {\it classical case} if $\Lambda$ is a Noetherian
complete local ring and $k$ is its residue field.
\end{definition}

\section{Notation and Conventions}
\label{section-notations-conventions}

\noindent
A ring is commutative with $1$. The  maximal ideal of a local ring $A$
is denoted by $\mathfrak{m}_A$. The set of positive integers is denoted
by $\mathbf{N} = \{1, 2, 3, \ldots\}$. If $U$ is an object of a
category $\mathcal{C}$, we denote by $\underline{U}$
the functor
$\Mor_\mathcal{C}(U, -): \mathcal{C} \to \textit{Sets}$, see
Remarks \href{https://stacks.math.columbia.edu/tag/06GK}{remarks cofibered groupoids (Tag 06GK)} (\href{https://stacks.math.columbia.edu/tag/06GQ}{item definition yoneda (Tag 06GQ)}).
Warning: this may conflict with the notation in other chapters where we
sometimes use $\underline{U}$ to denote $h_U(-) = \Mor_\mathcal{C}(-, U)$.

\medskip\noindent
Throughout this chapter $\Lambda$ is a Noetherian ring and
$\Lambda \to k$ is a finite ring map from $\Lambda$ to a field.
The kernel of this map is denoted $\mathfrak m_\Lambda$ and the
image $k' \subset k$. It turns out that $\mathfrak m_\Lambda$ is
a maximal ideal, $k' = \Lambda/\mathfrak m_\Lambda$ is a field, and
the extension $k/k'$ is finite. See discussion surrounding
(\href{https://stacks.math.columbia.edu/tag/06S2}{equation k prime (Tag 06S2)}).




\begin{definition}
\label{definition-completion-CLambda}
Let $\Lambda$ be a Noetherian ring and let $\Lambda \to k$ be a finite
ring map where $k$ is a field. We define {\it $\widehat{\mathcal{C}}_\Lambda$}
to be the category with
\begin{enumerate}
\item objects are pairs $(R, \varphi)$ where $R$ is a Noetherian complete
local $\Lambda$-algebra and where $\varphi : R/\mathfrak m_R \to k$ is a
$\Lambda$-algebra isomorphism, and
\item morphisms $f : (S, \psi) \to (R, \varphi)$ are local $\Lambda$-algebra
homomorphisms such that $\varphi \circ (f \bmod \mathfrak m) = \psi$.
\end{enumerate}
\end{definition}

\begin{definition}
\label{definition-smooth-morphism}
Let $\varphi : \mathcal{F} \to \mathcal{G}$ be a morphism of categories
cofibered in groupoids over $\mathcal{C}_\Lambda$.  We say  $\varphi$ is
{\it smooth} if it satisfies the following condition: Let $B \to A$ be
a surjective ring map in $\mathcal{C}_\Lambda$.  Let $y \in
\Ob(\mathcal{G}(B)), x \in \Ob(\mathcal{F}(A))$, and $y
\to \varphi(x)$ be a morphism lying over $B \to A$.  Then there
exists $x' \in \Ob(\mathcal{F}(B))$, a morphism $x' \to x$
lying over $B \to A$, and a morphism $\varphi(x') \to y$ lying
over $\text{id}: B \to B$, such that the diagram
$$
\xymatrix{
\varphi(x') \ar[r] \ar[dr] & y \ar[d] \\
& \varphi(x)
}
$$
commutes.
\end{definition}

\begin{definition}
\label{algebra-definition-formally-smooth}
Let $R \to S$ be a ring map.
We say $S$ is {\it formally smooth over $R$} if for every
commutative solid diagram
$$
\xymatrix{
S \ar[r] \ar@{-->}[rd] & A/I \\
R \ar[r] \ar[u] & A \ar[u]
}
$$
where $I \subset A$ is an ideal of square zero, a dotted
arrow exists which makes the diagram commute.
\end{definition}

\begin{lemma}
\label{lemma-smooth-morphism-power-series}
Let $R \to S$ be a ring map in $\widehat{\mathcal{C}}_\Lambda$. Then
the induced morphism
$\underline{S}|_{\mathcal{C}_\Lambda} \to \underline{R}|_{\mathcal{C}_\Lambda}$
is smooth if and only if $S$ is a power series ring over $R$.
\end{lemma}

\begin{proof}
Assume $S$ is a power series ring over $R$. Say $S = R[[x_1, \ldots, x_n]]$.
Smoothness of
$\underline{S}|_{\mathcal{C}_\Lambda} \to \underline{R}|_{\mathcal{C}_\Lambda}$
means the following (see Remark \href{https://stacks.math.columbia.edu/tag/06HJ}{remark compare smooth schlessinger (Tag 06HJ)}):
Given a surjective ring map $B \to A$ in
$\mathcal{C}_\Lambda$, a ring map $R \to B$, a ring map $S \to A$ such that
the solid diagram
$$
\xymatrix{
S \ar[r] \ar@{..>}[rd] & A \\
R \ar[u] \ar[r] & B \ar[u]
}
$$
is commutative then a dotted arrow exists making the diagram commute.
(Note the similarity with
Algebra, Definition \ref{algebra-definition-formally-smooth}.)
To construct the dotted arrow choose elements $b_i \in B$ whose images
in $A$ are equal to the images of $x_i$ in $A$. Note that
$b_i \in \mathfrak m_B$ as $x_i$ maps to an element of $\mathfrak m_A$.
Hence there is a unique $R$-algebra map $R[[x_1, \ldots, x_n]] \to B$
which maps $x_i$ to $b_i$ and which can serve as our dotted arrow.

\medskip\noindent
Conversely, assume
$\underline{S}|_{\mathcal{C}_\Lambda} \to \underline{R}|_{\mathcal{C}_\Lambda}$
is smooth. Let $x_1, \ldots, x_n \in S$ be elements whose images
form a basis in the relative cotangent space
$\mathfrak m_S/(\mathfrak m_R S + \mathfrak m_S^2)$ of $S$ over $R$.
Set $T = R[[X_1, \ldots, X_n]]$. Note that both
$$
S/(\mathfrak m_R S + \mathfrak m_S^2) \cong
R/\mathfrak m_R[x_1, \ldots, x_n]/(x_ix_j)
$$
and
$$
T/(\mathfrak m_R T + \mathfrak m_T^2) \cong
R/\mathfrak m_R[X_1, \ldots, X_n]/(X_iX_j).
$$
Let
$S/(\mathfrak m_R S + \mathfrak m_S^2) \to
T/(\mathfrak m_R T + \mathfrak m_T^2)$
be the local $R$-algebra isomorphism given by mapping
the class of $x_i$ to the class of $X_i$. Let
$f_1 : S \to T/(\mathfrak m_R T + \mathfrak m_T^2)$ be the
composition
$S \to S/(\mathfrak m_R S + \mathfrak m_S^2)
\to T/(\mathfrak m_R T + \mathfrak m_T^2)$.
The assumption that
$\underline{S}|_{\mathcal{C}_\Lambda} \to \underline{R}|_{\mathcal{C}_\Lambda}$
is smooth means we can lift $f_1$ to a map
$f_2 : S \to T/\mathfrak{m}_T^2$, then to a map
$f_3 : S \to T/\mathfrak{m}_T^3$, and so on, for all $n \geq 1$. Thus
we get an induced map $f : S \to T = \lim T/\mathfrak m_T^n$
of local $R$-algebras. By our choice of $f_1$, the map $f$ induces an
isomorphism
$\mathfrak m_S/(\mathfrak m_R S + \mathfrak m_S^2) \to
\mathfrak m_T/(\mathfrak m_R T + \mathfrak m_T^2)$
of relative cotangent spaces.
Hence $f$ is surjective by
Lemma \href{https://stacks.math.columbia.edu/tag/06GZ}{lemma surjective cotangent space (Tag 06GZ)}
(where we think of $f$ as a map in $\widehat{\mathcal{C}}_R$).
Choose preimages $y_i \in S$ of $X_i \in T$ under $f$. As $T$ is a
power series ring over $R$ there exists a local
$R$-algebra homomorphism $s : T \to S$ mapping $X_i$ to $y_i$.
By construction $f \circ s = \text{id}$. Then $s$ is injective.
But $s$ induces an isomorphism on relative cotangent spaces since
$f$ does, so it is also surjective by
Lemma \href{https://stacks.math.columbia.edu/tag/06GZ}{lemma surjective cotangent space (Tag 06GZ)}
again. Hence $s$ and $f$ are isomorphisms.
\end{proof}
\end{document}
